% Fuente: Auxiliar 4, Pregunta 2.
Una empresa generadora opera una central de ciclo combinado y debe decidir la
producción de energía durante un bloque horario mediante dos componentes de la
planta: la turbina a gas y el bloque de vapor asociado.

En una central de ciclo combinado, la turbina a gas genera electricidad y sus
gases de escape se aprovechan en un generador de vapor de recuperación para
producir vapor. Este vapor acciona una turbina de vapor, permitiendo generar
electricidad adicional. El costo de operación de la turbina a gas es de
$4\;[\$/\mathrm{MWh}]$, mientras que el costo de operación del bloque de vapor
es de $2\;[\$/\mathrm{MWh}]$.

La energía producida por ambos componentes se remunera al precio nodal de la
barra, que durante el bloque horario es igual a $5\;[\$/\mathrm{MWh}]$ para
ambas tecnologías. La empresa busca maximizar su beneficio operativo durante
dicho bloque.

Además, el coordinador del sistema exige mantener al menos $4\;[\mathrm{MW}]$
de generación en la turbina a gas para cumplir un requerimiento de reservas
operacionales. Por otra parte, debido a limitaciones en el suministro de
combustible, la turbina a gas no puede producir más de $16\;[\mathrm{MW}]$.
En total, la central puede producir hasta $20\;[\mathrm{MW}]$ entre ambos
componentes.

Para representar el acoplamiento físico del ciclo combinado, se considera que
la producción del bloque de vapor no puede exceder la mitad de la producción de
la turbina a gas.

\medskip
\noindent
\textbf{Nota:} Como se trata de un bloque horario de una hora, puede
identificarse numéricamente la potencia promedio $[\mathrm{MW}]$ con la energía
producida durante el bloque $[\mathrm{MWh}]$. En lo que sigue, modele el
problema usando variables en $[\mathrm{MW}]$.

\begin{enumerate}[label=\alph*)]
    \item Modele el problema como un problema de optimización lineal. Defina
    claramente las variables de decisión, la función objetivo y las
    restricciones. Luego, reformúlelo en forma estándar.

    \item Realice una representación gráfica del problema en el espacio de las
    variables de decisión originales y encuentre el óptimo gráficamente.
    Luego, a partir de la forma estándar, obtenga la solución básica factible
    óptima. Indique claramente cuáles son las variables básicas y no básicas.
    Finalmente, proponga un algoritmo que permita encontrar la solución óptima
    en el caso de un poliedro en $\mathbb{R}^n$, con $n>2$.

    \item Debido a una condición operativa adicional, la mitad de la energía
    producida en el bloque de vapor, más la energía producida por la turbina a
    gas, debe sumar al menos $17\;[\mathrm{MW}]$. Agregue esta nueva
    restricción al problema, actualizando la formulación de optimización.
    Determine la solución básica factible óptima del problema modificado.

    \item Debido a un aumento en el costo del agua desmineralizada y en los
    insumos de tratamiento térmico, el costo de operación del bloque de vapor
    sube a $6\;[\$/\mathrm{MWh}]$. Determine el nuevo punto óptimo revisando la
    base óptima obtenida en la parte c). Indique si la base cambia y, en caso
    afirmativo, identifique qué variable sale de la base y cuál entra.

    \item Debido a una restricción térmica adicional (límite de temperatura del
    condensador), el bloque de vapor ahora no puede producir más de
    $4.5\;[\mathrm{MW}]$. \textbf{Partiendo del problema de la parte c)}
    (con $c_V = 2\;[\$/\mathrm{MWh}]$), agregue esta nueva restricción y
    analice qué ocurre con la base previamente obtenida. Luego, encuentre la
    nueva solución básica factible óptima.
\end{enumerate}

